\section{Preliminares}

%\subsection*{Preliminaries}
\newcommand{\C}{\mathbb{C}}
\newcommand{\e}{\operatorname{exp}}
\newcommand{\re}{\operatorname{Re}}
\newcommand{\im}{\operatorname{Im}}

%\subsection*{Arithmetic and Geometry}
Sea \(\C := \R[i] \) la extensión de cuerpo de \(\R\), obtenida de adicionar \(i\) (\emph{unidad imaginaria}), elemento que satisface \(i^2 = -1\). Algunos detalles sobre \(\C\): 
\begin{itemize}[left = 0.3cm]
    \item Cualquier \(z \in \C\) es escrito como \( z =  x + i y  \), donde \( x\) e \(y \) en \(\R\) son la \emph{parte real} e \emph{imaginaria} de \(z\). Tambien, \(\overline{z}:= x -iy\) denota el \emph{conjugado complejo} de \(z\). 
    \[x = \operatorname{Re}(z) = \frac{z + \overline{z}}{2} \ \ \ \text{e} \ \ \ y = \operatorname{Im}(z)= \frac{z - \overline{z}}{2i}.\]
    \item \(\C \simeq \R^2\), por lo que tenemos interpretaciones geométricas naturales de las operaciones con números complejos.
    \item La función \((z,w) \mapsto \langle z, w\rangle := z \overline{w}\) define un producto hermitiano que induce la norma \(|z|^2 = z\overline{z} = x^2 + y^2\) (misma usual de \(\R^2\)). Desigualdades útiles:  
    \begin{enumerate}[label = \roman*., left = 1.2cm]
        \item \(|z + w| \leq |z| + |w|\). \hspace{2cm} ii. \(\displaystyle\frac{1}{z} = \frac{\overline{z}}{|z|^2} \) (en \(\C^\times\)). 
        \item[iii.] \(||z| - |w|| \leq |z-w|\).  \hspace{1.6cm} iv. \(|\re(z)|, |\im(z) |\leq |z|\).
    \end{enumerate}
    \item Podemos escribir \(z \in \C^\times\) en \emph{forma polar} como \(z = re^{i\theta} = r\operatorname{cis}(\theta)\), donde \(\theta = \operatorname{arg}(z) \) es el \emph{argumento} de \(z\) e \(r = |z| > 0 \) es su módulo. 
\end{itemize}

%\begin{definition}
%    A sequence \((z_n)_{n\in \N}\subset \C\) \emph{converges} to \(w \in \C\), denoted \(z_n \to w \), iff \(\displaystyle\lim_{n\to \infty} |z_n - w| = 0 \). Moreover \((z_n)\) is \emph{Cauchy} iff \(\forall \epsilon > 0, \exists N \in \N : \forall n,m\geq N, \ |z_n -z_m| < \epsilon \).   
%\end{definition}

\newcommand{\D}{\mathbb{D}}
Siguiendo notación estándar, dados \(z_0 \in \C\) e \(r>0\), \(B(z_0; r) := \{z \in \C : |z- z_0| < r\}, \ B[z_0; r] := \overline{B(z_0; r)}\) e \( \partial B(z_0; r) = \{z \in \C : |z-z_0| = r\}\). 
\begin{definition}
    \centering \(\D := \{z \in \C : |z| <1\}\). \hspace{1cm } e \hspace{1cm}\(\rotatebox{45}{\large $\boxdot$}\operatorname{diam}\Omega := \displaystyle\sup_{\Omega} |z -w|\).
\end{definition}
La topología es la inducida por la métrica, luego abiertos están caracterizados por bolas y vale todo resultado conocido, en particular relacionados a compacidad y conexidad. 

\newcommand{\comentario}[1]{\textcolor{gray}{\(\rightarrow \) #1}}
%\begin{theorem}
%    \textbf{1.1.} \(\C\) is complete (all Cauchy sequence converges to some \(w \in \C\)). \comentario{Follows from natural identification with \(\R^2 \)}
%\end{theorem}

%\begin{definition}
%    \(\Omega \subset \C\) is \emph{bounded} if \(\exists M > 0 : \forall z \in \Omega, \ |z| < M \ \leftrightharpoons \ \exists B(z_0;r) : \Omega \subset B(z_0;r)\). 
%\end{definition}

\subsection*{Funciones holomorfas}

\newcommand{\ho}{\mathcal{O}}
\begin{definition}
    Una función \(f: \Omega\ab \subset \C \to \C\) es \emph{holomorfa} en \(z_0 \in \Omega\) si el límite,  
    \[\- \hspace{2.5cm}f'(z_0):= \lim_{h\to 0} \frac{f(z_0 + h) -f(z_0)}{h}, \ \ \textcolor{red}{ \ \ \ \ \ !! \ h = h_1 + i h_2 \in \C}\]% \ \ \ \textcolor{red}{\ !!\  h \in \C }\]
    existe. Cuando existe, \(f'(z_0)\) es la  \emph{derivada} de \(f\) en \(z_0\). Decimos que es holomorfa en \(\Omega\), \(f \in \ho(\Omega)\), si \(\forall z_0 \in \Omega, \ \exists f'(z_0)\) y \emph{entera} si \(f \in \ho(\C)\). %\comentario{Remember that \(h \in \C  \)!!}
\end{definition}

\begin{example}
    \(\C[z]\subset \ho(\C)\) e \(\nicefrac{1}{z} \in \ho(\C^\times)\). Por otro lado, la conjugación \(z \mapsto \overline{z}\) \ NO es holomorfa. \comentario{Estudie el límite en los ejes principales}
\end{example}

\begin{note}
    Una reformulación equivalente y útil de ser holomorfa en \(z_0\) es que \(\exists \alpha \in \C : f(z_0 + h) = f(z_0)+ \alpha\cdot h  + o(h) \), donde \(\alpha = f'(z_0)\). 
\end{note}

\begin{proposition}
    \textbf{2.2.} Si \(f\) e \(g\) son holomorfas (en dominios apropiados), entonces \(f+g, \ fg, \ \nicefrac{f}{g}\) e \(g \circ f\) son holomorfas tambien. \comentario{Mismo negôcio de análisis en \(\R^n\)} 
    %\[(f+g)' = f' + g'; \ \ (fg)' = f'g + fg'; \ \ \left(\frac{f}{g}\right)' = \frac{f'g - fg'}{g^2}; \ \ ((g\circ f)(z))' \]
\end{proposition}

\alerta{
\begin{example}
    \(F: (x,y) \mapsto (x,-y)  \) es diff, mas \(F \sim  f : z \mapsto \overline{z}\) \ NO es holomorfa. 
\end{example}
}

\renewcommand{\hom}{\mathcal{L}}
%\begin{note}
%    Remember that \(F: \R^2 \to \R^2\) is diff in \(p\) iff \(\exists L \in \hom (\R^2, \R^2) : F(p+v) = F(p) + L\cdot v + o \|v\|\) as \(\|v\| \to 0\). When exists, \(L = JF(p)\) is unique.          
%\end{note}

\begin{definition}[\emph{Operadores de Riemann}]
    \(\displaystyle \partial_z := \nicefrac{1}{2} \ (\partial_x + \nicefrac{1}{i} \ \partial_y)\) \ e \ \(\displaystyle \partial_{\overline{z}} := \nicefrac{1}{2}\  (\partial_x - \nicefrac{1}{i} \  \partial_y)\).   
\end{definition}

\begin{proposition}
    \textbf{2.3.} Si \(f(z) = u(x,y) + iv(x,y)\) es holomorfa en \(z_0 = x_0 + i y_0\), entonces 
    \[f_{\overline{z}}(z_0) = 0 \text{ \ \ e \ \ } f'(z_0) = f_z(z_0) = 2u_z (z_0).\]
    Más aún, \(F(x,y) = (u(x,y), v(x,y))\) es diff e \(\det JF(x_0,y_0) = |f'(z_0)|^2\). 
\end{proposition}

\demo{Para la primera igualdad (ecuaciones de Cauchy-Riemann) note que,
\[\lim_{h_1 \to 0}\frac{f(z_0+ h_1)}{h_1} = f'(z_0) = \lim_{ih_2 \to 0} \frac{f(z_0 + ih_2)}{ih_2} \Rightarrow \begin{cases}
    u_x = v_y \\ 
    u_y = -v_x
\end{cases}.\] 
Las otras afirmaciones tambien son consecuencia de las ecuaciones de C-R, 
\[JF(x_0, y_0)(h)= \begin{bmatrix}
u_x & u_y \\ 
-u_y & u_x
\end{bmatrix} \begin{bmatrix}
h_1 \\ h_2
\end{bmatrix} = (u_x - i u_y)(h_1 + ih_2) = f'(z_0) \cdot h.\]
\vspace{-0.4cm}
%also \(\det JF(x_0,y_0) = u_x^2 + u_y^2 = |f'(z_0)|^2\).  
}

\begin{theorem}
    \textbf{2.4.} Sean \(\Omega \simeq U\ab\subset \R^2\), \(u,v \in C^1(U, \R^2)\) e \(f = u+ iv :  \Omega \to \C \). Si \(u\) e \( v\) satisfisfacen la ecuaciones de C-R, entonces \(f\in \ho(\Omega)\) e \(f'(z) = f_z\). \comentario{Consiste en escribir \(du \) e \(dv\) por extensión y reagrupar, véase \cite[pag. 13]{stein} }
\end{theorem}

\subsection*{Series de Potencias}

Los dos ejemplos más importantes son la serie de \(e^z\), cuya convergencia es absoluta en compactos, y la geométrica, que converge absolutamente en \(\D\).  
\begin{itemize}
    \item \(\displaystyle \exp(z) = \sum_{n=0}^\infty \frac{z^n}{n!}\). \comentario{\(\displaystyle \left|\sum \frac{z^n}{n!}\right| \leq \sum \frac{|z|^n}{n!} = \exp|z| < \infty\)} 
    \item \(\displaystyle \frac{1}{1-z} = \sum_{n=0}^\infty z^n, \ |z| <1\). \comentario{\(\displaystyle S_N = \frac{1-z}{1-z}\sum^N z^n = \frac{1- z^{N+1}}{1-z} \to \frac{1}{1-z}\)}
\end{itemize}
En general, una \emph{serie de potencias} es una suma de términos de la forma \( (a_nz^n)_{n\geq 0 }  \) y estamos interesados en su convergencia absoluta.  

\begin{theorem}
    \textbf{2.5.} Dada una serie de potencias \(\sum a_n z^n\), existe \(R \in [0, \infty]\) tal que:  
    \[\sum_{n=0}^\infty |a_n||z|^n = \begin{cases}
        L, &|z| < R \\ 
        \infty,  &|z|>R 
    \end{cases}.\]
    Más aún, si usamos la convención \(\nicefrac{1}{0} = \infty\) e \(\nicefrac{1}{\infty} = 0\), entonces \(R\) viene dado por la \emph{formúla de Hadamard}, \(\nicefrac{1}{R} = \lim \sup |a_n|^{\nicefrac{1}{n}}\). 
\end{theorem}

\demo{Sea \(L := \nicefrac{1}{R}\), para \(|z| < R\) escoja \(0< \epsilon \ll 1 : (L+\epsilon)|z|< r < 1\). Sigue que \(|a_n||z|^n \leq \{(L+ \epsilon)|z|\}^n = r^n\) (serie geometrica). Un argumento análogo deja ver que la serie diverge para \(|z| > R\).  }

%\alerta{
%%    \centering
\begin{note}
    El caso \(|z| = R\) no es concluyente (\textcolor{rojoscuro}{Ex. 1.19}).
\end{note}
%    \tcblower
%    \begin{example}
%        Estudie la convergencia en \(\partial \D\) de \(\sum nz^n, \ \sum \nicefrac{1}{n^2}\cdot z^n\) e \( \sum \nicefrac{1}{n} \cdot z^n \). 
%    \end{example}
%}
\ejemplo{
    Formúlas de Euler para seno e coseno \centering 
    \tcblower
     \- \hspace{0.5cm} {\scriptsize \(\bullet\)} \ \(\displaystyle \cos (z) = \sum^\infty_{n=0} (-1)^n\frac{z^{2n}}{(2n)!} = \frac{e^{iz}+ e^{-iz}}{2}\) \\ 
    \- \hfill {\scriptsize \(\bullet\)} \ \(\displaystyle \sin (z) = \sum_{n=0}^\infty (-1)^n \frac{z^{2n+1}}{(2n+1)!} = \frac{e^{iz} - e^{-iz}}{2i}\) \-\hspace{0.5cm}
}

\begin{theorem}
    \textbf{2.6.} \(f(z) = \sum a_n z^n \in \ho(B(0;R))\). La derivada tambien es una serie, 
    \[f'(z) = \sum_{n=0}^\infty na_nz^{n-1} \ \textcolor{gray}{=:g(z)} \text{ (término a término)  },\] 
    de mismo radio de convegencia de la serie original. \comentario{El radio es el mismo pues \(n^{\nicefrac{1}{n}} \to 1\) cuando \(n\to \infty\), lo demás no admite resumen, véase \cite[pag. 17]{stein}}
\end{theorem}

%\demo{
    %h, se trata de un argumento \(\nicefrac{\epsilon }{3}\), usando \vspace{-0.3cm}
    %\[f(z) = S_N(z) + E_N(z) = \sum_{n=0}^N a_nz^n + \sum_{N+1}^\infty a_nz^n.\] 
    %\vspace{-0.3cm}
%}

\begin{corollary}
    \textbf{2.7.} Una serie de potencias es infinitamente diferenciable es su disco de convergencia, y las derivadas son obtenidas por diferenciación término a término. 
\end{corollary}

\begin{note}
    Lo mismo vale para series centradas en \(z_0 \in \C^\times\), son simples translaciones. 
\end{note}

\subsection*{Integración en curvas}

\begin{definition}
    Una \emph{curva parametrizada suave} es una función \(z: [a,b] \to \C \) tal que \(\exists z'(t)\neq 0\) continua en \([a,b]\) e con límites laterales en \(t = a,b\). Denote, 
    \[P :=\{z \mid z:[a,b] \to \C \text{ es c. p. s.}, \ a<b \in \R\}.\]  
    Dos parametrizaciones \(z: [a,b] \to \C\) e \(\widetilde{z}:[c,d] \to \C\) son \emph{equivalentes}, \(z \sim \widetilde{z}\), si \(\exists ^1t :[a,b] \simeq [c,d]\) tal que \(t'(s)>0\) e \( \widetilde{z}(s) = z(t(s))\). 
    \[\star \ \ \nicefrac{P}{\sim} \ni [z] = \operatorname{Im}z =: \gamma \subset \C \ \leftarrow \text{ \emph{Curva suave} de orientación }z(a) \leadsto z(b).\]
\end{definition}

\begin{note}
    Es posible extender de manera natural estás definiciones a \emph{curvas suaves a trozos}. Nos interesan ahora en particular las \emph{curvas cerradas} \(z \mid z(a) = z(b)\).
\end{note}

\begin{example}
    \(\partial B(z_0; r) = z_0 + r e^{it}, \ t\in [0, 2\pi]\), orientación positiva \(\circlearrowleft\). 
\end{example}

\begin{definition}
    Dada una curva suave \(\gamma \subset \C\) parametrizada por \(z:[a,b] \to \C\) e \(f\) una función continua en \(\gamma\) definimos, 
    \[\int_\gamma f(z) \ dz := \int_a^b f(z(t))z'(t) \ dt.\]
\end{definition}

\begin{note}
    Es fácil ver que la integral no depende de la parametrización. \comentario{Planteé una parametrización equivalente e aplique regla de la cadena }
\end{note}

\begin{definition}
    \centering
    \(\displaystyle \operatorname{length}(\gamma):= \int_a^b |z'(t) | \ dt\). 
\end{definition}

\begin{proposition}
    \textbf{3.1.} La integración de funciones continuas sobre curvas satisface:% properties: 
    \vspace{-0.3cm}
    \begin{enumerate}[label = \roman*., left = 0.35cm]
        \item \(\displaystyle \int_\gamma f(z) + \alpha g(z) \ dz = \int_\gamma f(z) \ dz + \alpha \int_\gamma g(z) \ dz \). 
        \item \(\displaystyle \int_\gamma f(z) \ dz = - \int_{\gamma^-} f(z) \ dz\).\hfill iii. \ \(\displaystyle \left|\int_\gamma f(z) \ dz\right| \leq  \operatorname{length}(\gamma) \cdot \sup_{z \in \gamma } |f(z)|\). 
    \end{enumerate}
\end{proposition}

\hypertarget{thm33}{}

\begin{theorem}
    \textbf{3.2.} Sean \(F \in \ho(\Omega)\), \( \Omega \supset \gamma : w_1 \leadsto w_2\) suave e \(dF(z)=  f(z) \ dz\). Entonces, 
    \[\int_\gamma f(z) \ dz = \int_{\partial \gamma} F(z) = F(w_2) - F(w_1). \textcolor{gray}{\ \ \rightarrow  \blacksquare \ \text{(\emph{Stokes})}}\]
\end{theorem}

\vspace{-0.5cm}
\hypertarget{coro133}{}
\begin{corollary}
    \textbf{3.3.} En las hipótesis del \(\blacksquare\) \hyperlink{thm33}{3.2.}, si \(\gamma\) fuera cerrada, entonces \(\displaystyle \int_\gamma f(z) \ dz = 0 \).
\end{corollary}

\hypertarget{coro134}{}
\begin{corollary}
    \textbf{3.4.} Si \(\Omega\ab\subset \C\) conexo, \(f\in \ho(\Omega)\)  e \(f'\equiv 0\), entonces \(f\) es constante. \comentario{Basta tomar \(\gamma: w_0 \leadsto w\) suave e aplicar \(\blacksquare\) \hyperlink{thm33}{3.2.}}
\end{corollary}

